\chapter{Painful Sex (Dyspareunia)}

\section*{Definition}
Recurrent or persistent genital pain occurring before, during, or after sexual intercourse, classified as superficial (introital) or deep pelvic pain.

\section*{Pathophysiology}
Superficial dyspareunia involves vulvar, vestibular, or vaginal pathology (lack of lubrication, vulvodynia, vaginismus, infection, atrophy). Deep dyspareunia (pelvic pain with thrusting) indicates upper reproductive tract pathology (endometriosis, PID, ovarian cyst, adenomyosis) or non-gynecologic causes (IBS, interstitial cystitis).

\section*{Etiology}
\begin{itemize}
  \item Vaginal dryness (especially menopause or breastfeeding)
  \item Vaginismus (involuntary pelvic floor muscle spasm)
  \item Vulvodynia or vestibulodynia (chronic vulvar pain)
  \item Endometriosis
  \item Pelvic inflammatory disease (PID)
  \item Ovarian cysts
  \item Interstitial cystitis (painful bladder syndrome)
  \item Uterine fibroids or adenomyosis
  \item Irritable bowel syndrome
  \item Postpartum perineal changes or scar tissue
\end{itemize}

\section*{Indications for Evaluation}
Seek care if:
\begin{itemize}
  \item Severe or worsening symptoms
  \item Painful sex with fever, bleeding, discharge, or lasting beyond initial penetration attempts
  \item Accompanied by other concerning signs
\end{itemize}

\section*{Related Symptoms}
\begin{itemize}
  \item Vaginal dryness
  \item Pelvic pain
  \item Irregular bleeding
  \item Urinary symptoms
  \item Lower back pain
\end{itemize}
\textbf{Note:} Always consider serious underlying conditions when evaluating persistent painful sex (dyspareunia).

\section*{Self-Care / Home Management}
\begin{itemize}
  \item Rest and avoid activities that may aggravate the condition.
  \item Maintain adequate hydration and a balanced diet.
  \item Over-the-counter remedies may provide symptomatic relief (consult a pharmacist or doctor).
  \item Monitor symptoms and seek professional advice if they persist or worsen.
  \item Avoid self-medicating with prescription medications without medical guidance.
\end{itemize}

\section*{Diagnostic Approach}
\begin{itemize}
  \item A thorough history and physical examination are the first steps in evaluation.
  \item Laboratory tests (blood work, urinalysis) may be ordered based on clinical suspicion.
  \item Imaging studies (X-ray, ultrasound, CT, MRI) may be indicated when structural pathology is suspected.
  \item Specialist referral is arranged when the diagnosis remains uncertain or requires specific expertise.
\end{itemize}

\section*{Treatment / Management Overview}
\begin{itemize}
  \item Treatment targets the underlying cause identified through diagnostic evaluation.
  \item Supportive care and symptom management are provided while awaiting definitive therapy.
  \item Medications, lifestyle modifications, and physical therapies may be prescribed as appropriate.
  \item Follow-up ensures treatment effectiveness and allows timely adjustments to the care plan.
\end{itemize}

\clearpage